\section{Ambient Displays und Human-Computer Interaction} \todo{text}
% NEU: 2.7 Ambient Displays und HCI (Human-Computer Interaction)
% Dieses Kapitel verleiht deiner Arbeit ein sehr starkes wissenschaftliches Fundament. Nutze hierfür deine bereitgestellten wissenschaftlichen Paper:
% Grundlagen von Ambient Information Systems (AIS): Beschreibe Ambient Displays als Systeme, die nicht-kritische Informationen (wie das Wetter oder die Mondphase) kontinuierlich am Rand der menschlichen Wahrnehmung ("periphery of attention") darstellen
% . Im Gegensatz zu Smartphones oder Monitoren, die den vollen Fokus erfordern, können Ambient Displays "im Vorbeigehen" erfasst werden, ohne primäre Tätigkeiten zu unterbrechen
% .
% Calm Technology: Zitiere hier die extrem oft zitierten Pioniere Mark Weiser und John Seely Brown. Das Ziel von "Calm Computing" ist es, Technologie so zu gestalten, dass sie fließend zwischen dem Zentrum und der Peripherie unserer Aufmerksamkeit wechseln kann
% . Ein gutes Ambient Display informiert uns, ohne uns zu dominieren oder zu stressen
% .
% Praxisbeispiele in der HCI: Erwähne verwandte Forschungsansätze, bei denen Uhren als Ambient Displays genutzt wurden. Ein prominentes Beispiel aus der Literatur ist die Wherebouts Clock (die den Aufenthaltsort von Familienmitgliedern anzeigt) oder die Google Clock (die Termine durch farbliche Codierung im Hintergrund signalisiert)
% . Deine Wortuhr schlägt genau in diese Kerbe: Sie nutzt die Ästhetik eines Alltagsgegenstandes, um komplexe Smart-Home- und Wetterdaten unaufdringlich darzustellen
% .


% Für Kapitel 1.1: Motivation und Problemstellung
% (Hier begründen Sie die Relevanz Ihrer Arbeit über das Problem der Informationsüberflutung und den Lösungsansatz der "Calm Technology".)
% Textvorschlag:
% "In der heutigen Zeit sehen sich Nutzer durch die stetige Zunahme von vernetzten Geräten und Pervasive Computing mit einer wachsenden Informationsüberflutung (Information Overload) konfrontiert
% . Klassische Interfaces wie Smartphones oder PC-Bildschirme fordern meist die volle Aufmerksamkeit des Nutzers ein und reißen ihn aus seinem aktuellen Kontext heraus
% . Um dieser kognitiven Belastung entgegenzuwirken, forderten Mark Weiser und John Seely Brown bereits in den späten 1990er Jahren mit dem Paradigma der Calm Technology, dass Technologie nahtlos in den Hintergrund treten und Informationen unaufdringlich am Rande der Aufmerksamkeit – in der Peripherie – vermitteln sollte
% .
% Die Relevanz des vorliegenden Projekts liegt in der Anwendung genau dieses Prinzips auf den Smart-Home-Kontext. Während klassische Wortuhren primär ästhetischen Zwecken dienen und moderne Smart Displays häufig kognitiv überladen sind, zielt die hier entwickelte WordClock darauf ab, als sogenanntes Ambient Information System (AIS) zu fungieren
% . Durch die ästhetische, wortbasierte Anzeige relevanter Zusatzdaten (wie exakte Minuten, Wetter oder Mondphasen) wird eine Brücke zwischen Dekoration und multifunktionalem Smart-Home-Interface geschlagen
% . Die Informationen sind stets präsent, erfordern jedoch keine unmittelbare Interaktion und binden keine unnötigen kognitiven Ressourcen
% ."
% Für das NEUE Kapitel 2.2: Wissenschaftliches Umfeld: Ambient Displays und HCI
% (Hier schaffen Sie das theoretische Fundament und belegen, dass Ihr Ansatz state-of-the-art in der Forschung ist.)
% Textvorschlag:
% "Die konzeptionelle Basis für die erweiterte WordClock bilden Erkenntnisse aus dem Bereich der Mensch-Computer-Interaktion (HCI), speziell zu Ambient Information Systems und Peripheral Interaction.
% Ambient Information Systems (AIS) Ambient Information Systems präsentieren Informationen kontinuierlich und unaufdringlich im Hintergrund der menschlichen Wahrnehmung
% . Ihr Ziel ist es, den Nutzer über nicht-kritische Informationen (wie z. B. das Wetter) auf dem Laufenden zu halten, ohne ihn bei seinen primären Tätigkeiten zu stören
% . Diese Systeme nutzen die menschliche Fähigkeit zur präattentiven Verarbeitung, wodurch Informationen intuitiv und mit geringem mentalen Aufwand im Vorbeigehen erfasst werden können
% . Die WordClock greift dieses Prinzip auf: Die Anzeige der Zeit und des Wetters erfolgt visuell ansprechend und dezent, sodass sie die Wohnumgebung bereichert, anstatt als störendes technisches Gerät aufzufallen.
% Peripheral Interaction und Tangible Interfaces Ein weiteres relevantes Konzept ist die Peripheral Interaction. Im Alltag führen Menschen viele Handlungen am Rande ihrer Aufmerksamkeit aus (z. B. das Erfassen des Wetters beim Blick aus dem Fenster), können den Fokus aber jederzeit darauf richten, wenn es nötig ist
% . Die WordClock ermöglicht genau diesen fließenden Wechsel: Sie ruht in der Peripherie, kann aber bei gezielter Betrachtung detaillierte Auskunft geben
% . Zudem weist das System Eigenschaften eines Tangible User Interfaces (TUI) auf, da es digitalen Informationen (Wetter-APIs, Smart-Home-Daten) eine physische Form und Repräsentanz im Raum gibt, die über einen reinen Bildschirm hinausgeht
% .
% Smart Home und User Experience Im Bereich der Smart-Home-Forschung wird betont, dass intelligente Umgebungen hochgradig anpassbar sein müssen, um die Lebensqualität zu erhöhen
% . Eine effektive User Experience entsteht dann, wenn sich die Schnittstellen (Interfaces) an die Gewohnheiten der Nutzer anpassen und eine zentrale, vernetzte Kontrolle ermöglichen
% . Die Integration der WordClock in bestehende Smart-Home-Strukturen sowie die Möglichkeit der App-Steuerung tragen dieser Forderung nach personalisierbarer und effizienter Interaktion Rechnung."
% Für Kapitel 7: Diskussion der Ergebnisse
% (Ein Ausblick, wie Sie später bei der Auswertung Ihrer Ergebnisse auf die Theorie zurückgreifen können.)
% Textvorschlag (Ideen für später):
% "Bei der Diskussion der Ergebnisse in Kapitel 7 können Sie bewerten, inwiefern Ihr Prototyp die Anforderungen an ein Ambient Information System erfüllt hat. Sie können anhand der Kriterien für AIS diskutieren, ob die Anzeigeumschaltung (Zusatzinformationen wie Mondphasen) fließend genug ist, um nicht als ablenkend (distracting) empfunden zu werden
% . Sie können zudem argumentieren, dass Ihr System erfolgreich die Lücke schließt, indem es eine alltägliche Notwendigkeit (Zeitmessung) mit der Vision der Calm Technology
%  und den aktuellen Anforderungen an vernetzte Smart-Home-Umgebungen vereint
% ."
% Mit diesen Bausteinen haben Sie eine starke Brücke zwischen der praktischen Umsetzung (Ihrer WordClock) und der hochakademischen Forschung (HCI, Ubiquitous Computing, Ambient Displays) geschlagen. Genau das wertet eine ingenieurswissenschaftliche Arbeit für jeden Betreuer massiv auf!
% Lassen Sie mich wissen, ob wir einen der Aspekte (z.B. das Smart Home Protokoll-Kapitel) noch weiter mit diesen Quellen unterfüttern sollen!